% !TeX root = ../../../../../main.tex
% cd.tex

\subsection{CD（消費需要）曲線の導出}
\label{sec:chapter5_asad_cd}
CD曲線は、需要側の全方程式を集約したAD（生産需要）曲線を出発点として導出される。

AD曲線の導出過程（式\eqref{form:chapter5_asad_ad_p_H_y_H_modified_eq}TODO）において、名目GDP（\(p_t^Hy_t^H\)）は以下の式によって決定されることが示された。
\begin{equation}
\label{eq:cd_ad_curve_gdp}
    p_t^Hy_t^H=\frac{\Gamma_t}{\beta_t^H(1+i_t^H)\operatorname{E}_t[\lambda_{t+1}^H]}
\end{equation}
ここで\(\Gamma_t\equiv\alpha^H+(1-\alpha^F)\operatorname{E}_t[\mathcal{B}_t\Lambda_{\infty}]\)は、自国政策から独立したパラメータと期待項のまとまりである。

一方で、財市場の均衡条件や一物一価の法則、名目為替レートの決定式を組み合わせることで、名目GDPと名目総消費（\(p_t^{H\to W}c_t^{H\to W}\)）の間には以下の関係が成り立つことも示されている（式\eqref{form:chapter5_asad_ad_p_H_y_H_eq}TODO）。
\begin{equation}
\label{eq:cd_nominal_conversion}
    p_t^Hy_t^H=\Gamma_tp_t^{H\to W}c_t^{H\to W}
\end{equation}

これら2つの式は、生産平面における需要（AD曲線）と、消費平面における需要（名目総消費）が、交易条件等の係数\(\Gamma_t\)を介して表裏一体であることを示している。
式\eqref{eq:cd_nominal_conversion}を式\eqref{eq:cd_ad_curve_gdp}の左辺に代入すると以下のようになる。
\begin{equation}
\label{eq:cd_substitution}
    \Gamma_tp_t^{H\to W}c_t^{H\to W}=\frac{\Gamma_t}{\beta_t^H(1+i_t^H)\operatorname{E}_t[\lambda_{t+1}^H]}
\end{equation}

両辺の\(\Gamma_t\)が相殺されることにより、生産側の複雑な波及経路が消去され、消費者物価指数\(p_t^{H\to W}\)と実質消費\(c_t^{H\to W}\)の純粋な需要関係を示す以下の\textbf{CD曲線}が導出される。
\begin{equation}
\label{eq:cd_curve}
    p_t^{H\to W}=\frac{1}{c_t^{H\to W}}\times\frac{1}{\beta_t^H(1+i_t^H)\operatorname{E}_t[\lambda_{t+1}^H]}
\end{equation}

形状:式\eqref{eq:cd_curve}からわかるように、CD曲線は\(p_t^{H\to W}\)と\(c_t^{H\to W}\)が反比例の関係にあるため、消費平面\((c_t^{H\to W},p_t^{H\to W})\)において右下がりの直角双曲線として描かれる。

シフト要因:曲線の位置は、名目利子率\(i_t^H\)や所得の期待限界効用\(\operatorname{E}_t[\lambda_{t+1}^H]\)によって決定される。中央銀行が名目利子率\(i_t^H\)を引き下げるか、あるいは将来の期待限界効用\(\operatorname{E}_t[\lambda_{t+1}^H]\)を低下させる（将来の消費が潤沢であると約束する）と、分母が小さくなるためCD曲線は右上（外側）へとシフトし、消費需要を拡大させる。